I used \href{https://github.com/hathach/tinyusb}{TinyUSB} to generate a virtual COM port for data transfer to a PC as it support the U575 and had open-source code for implementation. Before setting up the code, I needed to configure the STM into device-only full-speed and ensure the clock configuration had the USB peripheral at 48 MHz exactly.
\nl
Implementing TinyUSB was a matter of configuring \verb|tusb_config.h| and \verb|tusb_optn.h| to reflect the required settings like setting the OS as ThreadX, the MCU as the STM32U5, setting the class to CDC, etc. By default, the RX and TX buffer size \verb|CFG_TUD_CDC_TX_BUFSIZE| is 64 in full-speed mode, which is far too small for image transfer. Initial testing showed that this corresponded to roughly one row of pixels per few seconds. By increasing these to 1024+, data transfer speeds increased greatly.
\nl
To differentiate the different types of data being received (CSA vs. image data), a basic packet structure was defined:

\begin{itemize}
    \item \verb|Byte [1:0]|: Packet start sequence (\texttt{0xAA55})
    \item \verb|Byte [2]|: Packet ID (\texttt{0x01} = Image Data, \texttt{0x02} = CSA Data)
    \item \verb|Byte [6:3]|: Data length ($N$ bytes)
    \item \verb|Byte [6+N-1 : 7]|: Data
    \item \verb|Byte [6+N]|: Packet end sequence (\texttt{0xBB})
\end{itemize}

As the data length is defined, there is no risk for a premature start/end during nominal operations.
\nl
I used \href{https://www.der-hammer.info/pages/terminal.html}{HTerm} to verify behaviour during initial testing as it allowed me to view all data received in an intuitive format, and used \href{https://gemini.google.com/app}{Gemini (Gen AI)} to generate a dashboard to receive image data (see \autoref{fig:v1-dashboard}). The dashboard presents an easy way to configure the image sensor by using drop-downs and checkboxes to mix-and-match modes, with information icons to explain how certain features work, and allows you to save all received images to a set folder on your computer.

\begin{figure}
    \centering
    \includegraphics[width=0.7\textwidth]{Images/v1/Software/dashboard.png}
    \caption{Test pattern displayed on dashboard generated by Gemini (Gen AI).}
    \label{fig:v1-dashboard}
\end{figure}

Another basic packaet structure was defined for sending commands from the dashboard to C2Sv1.0:

\begin{itemize}
    \item \verb|Byte [0]|: Start condition (0xFF)
    \item \verb|Byte [1]|: Colour mode (0=RGB565, 1=Grayscale, 2=YCbCr)
    \item \verb|Byte [2]|: Resolution (0=VGA, 1=QVGA)
    \item \verb|Byte [3]|: Test Pattern (0=Disabled, 1=Enabled)
    \item \verb|Byte [4]|: End condition (0xFE)
\end{itemize}

C2Sv1.0 listens on the VCOM port for incoming data, and when the start and end conditions are met, the image capture process is begun by initialising the image sensor in the desired mode.